%% LyX 2.4.4 created this file.  For more info, see https://www.lyx.org/.
%% Do not edit unless you really know what you are doing.
\documentclass[english]{article}
\usepackage[T1]{fontenc}
\usepackage[latin9]{inputenc}
\usepackage{units}
\usepackage{enumitem}
\usepackage{amsmath}
\usepackage{amsthm}
\usepackage{amssymb}
\usepackage{geometry}
\geometry{verbose,tmargin=1in,bmargin=1in,lmargin=1in,rmargin=1in}

\makeatletter
%%%%%%%%%%%%%%%%%%%%%%%%%%%%%% Textclass specific LaTeX commands.
\numberwithin{equation}{section}
\numberwithin{figure}{section}
\newlength{\lyxlabelwidth}      % auxiliary length 

%%%%%%%%%%%%%%%%%%%%%%%%%%%%%% User specified LaTeX commands.
\usepackage{fancyhdr}
\usepackage{lastpage}
\pagestyle{fancy}
\usepackage{enumitem}
\usepackage{ifthen}
\everymath=\expandafter{\the\everymath\displaystyle}
%\DeclareMathSizes{10}{10}{10}{10}
\usepackage{multicol}

\usepackage{tikz}
\usetikzlibrary{patterns}
\usetikzlibrary{plotmarks}
\usepackage{pgfplots}
\pgfplotsset{compat=newest}

\usepackage{calc}
\usepackage[nomessages]{fp}

\usepackage{datetime}
% Added by lyx2lyx
\setlength{\parskip}{\smallskipamount}
\setlength{\parindent}{0pt}

\makeatother

\usepackage{babel}
\begin{document}
\renewcommand{\author}{Dr.\,James Ashe}

\newcommand{\classNum}{336}

\newcommand{\classSec}{A}

\newcommand{\examType}{Quiz 1.1-1.2}

\renewcommand{\date}{\formatdate{30}{8}{2013}}

%%First page's header%%%%%%%%%%%%%%%%%%%%%%
\fancypagestyle{firststyle} %%header settings for first pagr
{
	\fancyhead[L]{MTH \classNum{}(\classSec{}) \author{}\\
	\textbf{\examType{}} ~\date{}}
	\fancyhead[C]{}
	\fancyhead[R]{\textbf{Name:\hspace{4cm}}}
	\setlength{\headsep}{14pt}
}
%%%%%%%%%%%%%%%%%%%%%%%%%%%%%%%%%%%%%%%%%%

%%Next page's header%%%%%%%%%%%%%%%%%%%%%%%
\setlist{nolistsep}
\lhead{\classNum{}(\classSec{})} 
\chead{\textbf{Name: Solutions}} 
\cfoot{\thepage{}~of~\pageref{LastPage}} 
\setlength{\headsep}{5pt} 
%%%%%%%%%%%%%%%%%%%%%%%%%%%%%%%%%%%%%%%%%%%

%%Various settings; see the preamble for other settings%%%%%
%\setenumerate[0]{leftmargin=12pt,itemindent=0pt,labelwidth=0pt}
\setlist{topsep=.5pt}
\renewcommand{\labelenumi}{\textbf{\arabic{enumi}.}}
\renewcommand{\labelenumii}{\textbf{\alph{enumii}.}}
\thispagestyle{firststyle}
%%%%%%%%%%%%%%%%%%%%%%%%%%%%%%%%%%%%%%%%%%

\newcommand{\xmax}{0}
\newcommand{\xmin}{-\xmax}
\newcommand{\xstep}{0}
\newcommand{\ymax}{0}
\newcommand{\ymin}{-\ymax}
\newcommand{\ystep}{0} 
\newcommand{\dspace}{\vspace{1cm}}
\begin{enumerate}[resume]
\item Let $\mathbf{x}=(2,-5)$ and $\mathbf{y}=(0,3)$. Complete the following:\begin{multicols}{2}

\begin{enumerate}[resume]
\item $\mathbf{x}+\mathbf{y}=(2+0,-5+3)=(2,-2)$\vfill
\item Sketch $\mathbf{x}+\mathbf{y}$. 
\item Sketch $\mathbf{x}-\mathbf{y}$. $(0,-8)$
\item Sketch $\frac{1}{2}\mathbf{x}$. $(1,-\nicefrac{5}{2})$\columnbreak
\item []\renewcommand{\xmax}{3}
\renewcommand{\xmin}{-1}
\renewcommand{\xstep}{0} %must be xmin+n
\renewcommand{\ymax}{4}
\renewcommand{\ymin}{-6}
\renewcommand{\ystep}{0} %must be ymin+n
%%%%%%%
\begin{tikzpicture} 
\begin{axis}[height=5cm, 
	width=5cm, 
	axis lines=middle,
	minor x tick num={1},
	minor y tick num={1},
	grid=both,
	%xtick={0,50,...,250},
	%ytick={0,10,20,...,40},
	%axis line style={-latex}, 
	xmin=\xmin,xmax=\xmax, 
	ymin=\ymin,ymax=\ymax,
	%restrict y to domain=-1:10,
	%no markers, 
	xlabel={$$},ylabel={$$}, 
	%every axis x label/.append style={anchor=north}, 
	%every axis y label/.append style={anchor=east}
	]
	\addplot[->,color=black,very thick,solid] coordinates {(0,0) (2,-5)} node [left] {\textbf{x}};
	\addplot[->,color=black,very thick,solid] coordinates {(0,0) (0,3)} node [left] {\textbf{y}};
	\addplot[->,color=green,thick,solid] coordinates {(0,0) (2,-2)} node [right] {\textbf{a.}};
	\addplot[->,color=blue,thick,solid] coordinates {(0,0) (1,-2.5)} node [left] {\textbf{c.}};
	\addplot[->,color=red,thick,solid] coordinates {(0,3) (2,-5)} node [right] {\textbf{b.}};
	%\draw[blue,thick,->] (axis cs:0,0) -- (axis cs:-1,1) node [left] {y};
	%%\addplot[domain=\xmin:\xmax,thick,samples=100,draw=red,<->]{-(x-2)^2+5} [yshift=8pt] node[pos=0.7,right] {$g(x)$};
	
\end{axis}
\end{tikzpicture}
\end{enumerate}
\end{multicols}
\item Determine if the point $(-6,2)$ is on the line defined parametrically
by $\mathbf{x}=(0,1)+t(-3,1)$.

\begin{enumerate}[resume]
\item []$(0,1)+t(-3,1)=(-6,2)\Rightarrow0+-3t=-6\Rightarrow t=2$.\\
So then $(0,1)+2(-3,1)=(-6,3)\neq(-6,2)$. So the point is \emph{not
}on the line.\vfill
\end{enumerate}
\item Let $(1,1,1)$ and $(2,0,1)$ be two points on a plane $H$. Find
a line on the plane. 

\begin{enumerate}[resume]
\item []$(1,1,1)-(2,0,1)=(-1,1,0)$ is the direction vector of the line
connecting the two points. So then $\ell=(1,1,1)+t(-1,1,0)$ for $t\in\mathbb{R}$
is on the plane (answers may vary).\vfill
\end{enumerate}
\item Let $\mathbf{x}$ and $\mathbf{y}$ be two vectors such that $\mathbf{x}\cdot\mathbf{y}=0$.

\begin{enumerate}[resume]
\item What is the angle between $\mathbf{x}$ and \textbf{$\mathbf{y}$}?\textbf{
}

\begin{enumerate}[resume]
\item []The dot product is 0 meaning that the vectors are orthogonal\textbf{.
}So the angle is \textbf{$\frac{\pi}{2}$ }or $90^{\circ}$. 
\end{enumerate}
\item If $\mathbf{x}=(1,1)$ find a vector parallel to $\mathbf{y}$.

\begin{enumerate}[resume]
\item []$(1,1)\cdot(y_{1},y_{2})=1y_{1}+1y_{2}=0\Rightarrow y_{1}=-y_{2}$.
So $(1,-1)$ is parallel to $\mathbf{y}$ (answers may vary). \\
Note that and vector parallel to $\mathbf{y}$ will be equal to $c\mathbf{y}$
for some scalar $c\in\mathbb{R}$.\vfill
\end{enumerate}
\end{enumerate}
\end{enumerate}

\end{document}
